
\documentclass[12pt]{article}
\usepackage{amsmath, amsthm, amsfonts, mathtools, xcolor, caption}

\numberwithin{equation}{section}

\usepackage[style=numeric,sorting=none,giveninits=true]{biblatex}
\addbibresource{bibliography.bib}

\usepackage{etoolbox}
\usepackage[margin=1in]{geometry}

\DeclareMathAlphabet{\mymathbb}{U}{BOONDOX-ds}{m}{n}

\usepackage{comment}

%%%%%%%%%%%%%%%%%%%%%%%%%%%%%%%%%%%%%%%%%
%       Theorem Environments            %
%%%%%%%%%%%%%%%%%%%%%%%%%%%%%%%%%%%%%%%%%

\newtheorem{theorem}{Theorem}[section]
\newtheorem{lemma}[theorem]{Lemma}
\newtheorem{corollary}[theorem]{Corollary}
\newtheorem{proposition}[theorem]{Proposition}
\theoremstyle{definition}
\newtheorem{remark}{Remark}
\newtheorem{definition}{Definition}
\newtheorem{notation}{Notation}
\newtheorem{example}{Example}
\newtheorem{conjecture}[theorem]{Conjecture}
\newtheorem{openproblem}[theorem]{Open Problem}

%%%%%%%%%%%%%%%%%%%%%%%%%%%%%%%%%%%%%%%%%
%               New Commands            %
%%%%%%%%%%%%%%%%%%%%%%%%%%%%%%%%%%%%%%%%%

\newcommand{\ev}[1]{( #1 )} 
\newcommand{\ew}[1]{\left( #1 \right)} 

\newcommand{\N}{\mathbb{N}}
\newcommand{\Z}{\mathbb{Z}}
\newcommand{\C}{\mathbb{C}}
\newcommand{\Q}{\mathbb{Q}}
\newcommand{\R}{\mathbb{R}}

%%%%%%%%%%%%%%%%%%%%%%%%%%%%%%%%%%%%%%%%%
%               Document                %
%%%%%%%%%%%%%%%%%%%%%%%%%%%%%%%%%%%%%%%%%

\let\underbrace\LaTeXunderbrace
\let\overbrace\LaTeXoverbrace


\begin{document}

\section*{Challenge 1: The Hadwiger Conjecture}
\refstepcounter{section}

\begin{definition}[Hadwiger Number]
Given a graph $G$, the \textit{Hadwiger number} of $G$ is the largest integer $t$ such that the complete graph $K_t$ on $t$ vertices is a minor of $G$.
We denote the Hadwiger number by $h\ev{G}$.
\end{definition}


\begin{conjecture}[Hadwiger's Conjecture
\mbox{\cite[Conjecture 1]{steiner}}]\label{hadw}
Let $r\in \mathbb{N}$.
For every finite simple graph $G$, we have:
$$
r \leq \chi\ev{G} \ \Longrightarrow \  r \leq h\ev{G}
$$
where $\chi\ev{G}$ denotes the chromatic number of $G$.
\end{conjecture}

While Conjecture \ref{hadw} is not difficult for values of $r\leq 3$, already for $r=4$ it is equivalent to the 4-colour theorem. 
Robertson, Seymour and Thomas proved the above for $r= 5$ \cite{RobertsonSeymourThomas1993Hadwiger}. 

%%%%%%%%%%%%%%%%%%%%%%%%%%%%%%%%%%%%%%%%%%%%%%%%%%
\printbibliography
%%%%%%%%%%%%%%%%%%%%%%%%%%%%%%%%%%%%%%%%%%%%%%%%%%
\end{document}
